\documentclass[10pt,a4paper]{article}
\usepackage[a4paper,margin=14mm,top=12mm,bottom=14mm]{geometry}
\usepackage{fontspec}
\setmainfont{Liberation Sans}
\usepackage{polyglossia}
\setdefaultlanguage{french}
\usepackage{amsmath}
\usepackage{xcolor}
\usepackage{graphicx}
\usepackage[most]{tcolorbox}
\usepackage{tabularx,booktabs,array,multirow}
\usepackage{enumitem}
\usepackage{fancyhdr}
\usepackage{lastpage}
\usepackage[hidelinks]{hyperref}
\definecolor{uimm}{HTML}{005677}
\definecolor{cyan}{HTML}{13A2DD}
\definecolor{fond}{HTML}{E6F2FA}
\definecolor{ok}{HTML}{2E7D32}
\definecolor{ko}{HTML}{B3261E}
\definecolor{violet}{HTML}{7B1FA2}
\setlength{\parindent}{0pt}
\setlength{\parskip}{3pt}
\setlist{nosep}
\setlist[itemize]{leftmargin=13pt,label=\textcolor{cyan}{\textbullet}}
\newcommand{\titre}[1]{\vspace{5pt}{\color{uimm}\large\bfseries #1}\par\vspace{-3pt}{\color{cyan}\rule{\linewidth}{1.2pt}}\par\vspace{0pt}}
\tcbset{boite/.style={colback=fond,colframe=uimm,boxrule=0.6pt,arc=2pt,left=6pt,right=6pt,top=4pt,bottom=4pt},
        defi/.style={colback=violet!6,colframe=violet,boxrule=0.8pt,arc=2pt,left=6pt,right=6pt,top=4pt,bottom=4pt,
                     title={\bfseries Défi},coltitle=white,colbacktitle=violet,fonttitle=\small}}
\newcommand{\bandeau}[3]{%
\begin{tcolorbox}[colback=uimm,colframe=uimm,arc=2pt,left=4pt,right=8pt,top=4pt,bottom=4pt]
\begin{minipage}[c]{18mm}\setlength{\fboxsep}{1.2pt}\fcolorbox{white}{white}{\includegraphics[height=16.5mm]{logo_uimm.jpg}}\end{minipage}\hspace{4mm}%
\begin{minipage}[c]{\dimexpr\linewidth-22mm\relax}\color{white}
{\small #1}\\[1pt]{\LARGE\bfseries #2}\\[2pt]{\large #3}\hfill{\small CQPM Technicien qualité}
\end{minipage}
\end{tcolorbox}}
\newcommand{\reponse}[1][1]{\par\vspace{2pt}{\color{gray!60}\foreach \i in {1,...,#1}{\rule{\linewidth}{0.3pt}\par\vspace{9pt}}}}
\usepackage{pgffor}
\pagestyle{fancy}\fancyhf{}
\renewcommand{\headrulewidth}{0pt}
\newcommand{\pied}[1]{\fancyfoot[L]{\footnotesize\color{uimm}#1}\fancyfoot[R]{\footnotesize\color{uimm}Page \thepage{} sur \pageref{LastPage}}}

\usepackage{tikz}
\usepackage{longtable}
\renewcommand{\bandeau}[3]{%
\begin{tcolorbox}[colback=uimm,colframe=uimm,arc=2pt,left=4pt,right=8pt,top=4pt,bottom=4pt]
\begin{minipage}[c]{18mm}\setlength{\fboxsep}{1.2pt}\fcolorbox{white}{white}{\includegraphics[height=16.5mm]{logo_uimm.jpg}}\end{minipage}\hspace{4mm}%
\begin{minipage}[c]{\dimexpr\linewidth-22mm\relax}\color{white}
{\small #1}\\[1pt]{\LARGE\bfseries #2}\\[2pt]{\large #3}
\end{minipage}
\end{tcolorbox}}
% en-tête d'une fiche : code, titre, question, bloc, durée, organisation, niveau
\newcommand{\entete}[7]{%
\begin{tcolorbox}[colback=uimm,colframe=uimm,arc=2pt,left=4pt,right=6pt,top=4pt,bottom=4pt]
\begin{minipage}[c]{18mm}\setlength{\fboxsep}{1.2pt}\fcolorbox{white}{white}{\includegraphics[height=16.5mm]{logo_uimm.jpg}}\end{minipage}\hspace{3mm}%
\begin{minipage}[c]{\dimexpr\linewidth-58mm\relax}\color{white}\raggedright
{\small Base SPC · #4}\\[1pt]{\Large\bfseries #1 · #2}\\[2pt]{\itshape #3}
\end{minipage}\hfill
\begin{minipage}[c]{33mm}\color{white}\small\raggedleft
Durée : \textbf{#5}\\ \textbf{#6}\\ Niveau #7 sur 3
\end{minipage}
\end{tcolorbox}}
\newcommand{\tagE}{{\setlength{\fboxsep}{1.5pt}\colorbox{cyan}{\color{white}\scriptsize\bfseries Ellistat}}}
\newcommand{\tagW}{{\setlength{\fboxsep}{1.5pt}\colorbox{uimm}{\color{white}\scriptsize\bfseries Page web}}}
\tcbset{pred/.style={colback=orange!6,colframe=orange!80!black,boxrule=0.6pt,arc=2pt,left=6pt,right=6pt,top=3pt,bottom=1pt,
                    title={\bfseries Je prédis, avant de commencer},coltitle=white,colbacktitle=orange!80!black,fonttitle=\small},
        retenir/.style={colback=fond,colframe=uimm,boxrule=0.6pt,arc=2pt,left=6pt,right=6pt,top=3pt,bottom=3pt,
                    title={\bfseries À retenir},coltitle=white,colbacktitle=uimm,fonttitle=\small},
        plus/.style={colback=violet!6,colframe=violet,boxrule=0.6pt,arc=2pt,left=6pt,right=6pt,top=3pt,bottom=3pt,
                    title={\bfseries Pour aller plus loin},coltitle=white,colbacktitle=violet,fonttitle=\small},
        act/.style={colback=uimm!10,colframe=uimm,boxrule=0.5pt,arc=2pt,left=5pt,right=5pt,top=2pt,bottom=2pt,before skip=7pt,after skip=3pt},
        verif/.style={colback=orange!6,colframe=orange!80!black,boxrule=0.6pt,arc=2pt,left=6pt,right=6pt,top=3pt,bottom=3pt,
                    title={\bfseries À vérifier sur poste avant la première séance},coltitle=white,colbacktitle=orange!80!black,fonttitle=\small}}
\newcommand{\fichier}[1]{\texttt{\small #1}}
\newcommand{\menu}[1]{\textsf{\textbf{#1}}}
\newcommand{\champ}[1]{\rule{#1}{0.3pt}}

\pied{Base SPC · fiche A4 · Le procédé se dérègle}
\begin{document}
\thispagestyle{fancy}
\entete{A4}{Le procédé se dérègle}{Ma carte voit-elle le changement, et au bout de combien de lots ?}{Cartes aux attributs}{30 min}{en binôme}{2}
\begin{tcolorbox}[boite]\textbf{Le but.} Surveiller une production avec des limites figées et repérer un dérèglement.\par\textbf{Il vous faut.} La page « Tirer un lot », scénario donné par le formateur · La carte papier de A2 (ou une nouvelle)\end{tcolorbox}
\begin{tcolorbox}[pred]Si le taux de défectueuses double, la carte sonnera-t-elle dès le premier lot ?\reponse\end{tcolorbox}
\titre{À faire}
\begin{enumerate}[label=\textcolor{cyan}{\textbf{\arabic*.}},leftmargin=15pt,itemsep=1.5pt]
\item \tagW\ Page « Tirer un lot », nouveau scénario donné par le formateur, même taille de lot (50). Tirer 20 lots et calculer les limites comme en A2 : elles sont maintenant \textbf{figées}.
\item Continuer : tirer les lots 21 à 35 un par un, placer chaque point, sans recalculer les limites.
\item Noter le numéro du premier lot qui sort des limites.
\item \tagE\ Coller les 35 lots dans Ellistat. Carte p avec la \menu{Phase I} du lot 1 au lot 20.
\item \tagE\ Régler ensuite la Phase I sur les 35 lots. Regarder ce que deviennent les limites.
\end{enumerate}
\vspace{3pt}{\renewcommand{\arraystretch}{1.35}\small\begin{tabularx}{\linewidth}{|l|*{12}{>{\centering\arraybackslash}X|}}\hline
\textbf{Lot} & 1 & 2 & 3 & 4 & 5 & 6 & 7 & 8 & 9 & 10 & 11 & 12\\\hline
\textbf{Défectueuses} &  &  &  &  &  &  &  &  &  &  &  & \\\hline\end{tabularx}}\par\vspace{2pt}
{\renewcommand{\arraystretch}{1.35}\small\begin{tabularx}{\linewidth}{|l|*{12}{>{\centering\arraybackslash}X|}}\hline
\textbf{Lot} & 13 & 14 & 15 & 16 & 17 & 18 & 19 & 20 & 21 & 22 & 23 & 24\\\hline
\textbf{Défectueuses} &  &  &  &  &  &  &  &  &  &  &  & \\\hline\end{tabularx}}\par\vspace{2pt}
{\renewcommand{\arraystretch}{1.35}\small\begin{tabularx}{\linewidth}{|l|*{12}{>{\centering\arraybackslash}X|}}\hline
\textbf{Lot} & 25 & 26 & 27 & 28 & 29 & 30 & 31 & 32 & 33 & 34 & 35 & \\\hline
\textbf{Défectueuses} &  &  &  &  &  &  &  &  &  &  &  & \\\hline\end{tabularx}}\par\vspace{2pt}
\begin{center}\begin{tikzpicture}[x=0.446cm,y=0.31cm]\draw[gray!45,very thin] (0,0) grid (35,12);\draw[thick] (0,0)--(35,0);\draw[thick] (0,0)--(0,12);\node[left,font=\scriptsize] at (0,0) {0};\node[left,font=\scriptsize] at (0,2) {2};\node[left,font=\scriptsize] at (0,4) {4};\node[left,font=\scriptsize] at (0,6) {6};\node[left,font=\scriptsize] at (0,8) {8};\node[left,font=\scriptsize] at (0,10) {10};\node[left,font=\scriptsize] at (0,12) {12};\node[below,font=\scriptsize] at (1,0) {1};\node[below,font=\scriptsize] at (5,0) {5};\node[below,font=\scriptsize] at (10,0) {10};\node[below,font=\scriptsize] at (15,0) {15};\node[below,font=\scriptsize] at (20,0) {20};\node[below,font=\scriptsize] at (25,0) {25};\node[below,font=\scriptsize] at (30,0) {30};\node[below,font=\scriptsize] at (35,0) {35};\node[right,font=\scriptsize] at (35,0) {0 \%};\node[right,font=\scriptsize] at (35,2) {4 \%};\node[right,font=\scriptsize] at (35,4) {8 \%};\node[right,font=\scriptsize] at (35,6) {12 \%};\node[right,font=\scriptsize] at (35,8) {16 \%};\node[right,font=\scriptsize] at (35,10) {20 \%};\node[right,font=\scriptsize] at (35,12) {24 \%};\node[rotate=90,font=\scriptsize] at (-2.47,6.0) {pièces défectueuses};\node[below,font=\scriptsize] at (17.5,-1.3) {lot};\end{tikzpicture}\end{center}
\newpage
\titre{J'observe, je décide}
\begin{enumerate}[label=\textcolor{cyan}{\textbf{Q\arabic*.}},leftmargin=22pt]
\item Premier lot hors limites : n° \champ{10mm}. Combien de lots se sont écoulés depuis le lot 21 ?\reponse[1]
\item Avant ce lot, voyait-on déjà quelque chose sur la carte ? Quoi ?\reponse[2]
\item Avec les limites calculées sur les 35 lots, combien de points sortent ? Pourquoi ?\reponse[2]
\item Demandez au formateur de révéler le scénario. Votre carte avait-elle raison ?\reponse[1]
\end{enumerate}
\begin{tcolorbox}[retenir]On calcule les limites sur une période stable, puis on les fige. Recalculer avec des lots déréglés élargit les limites et cache le problème.\end{tcolorbox}
\begin{tcolorbox}[plus]Variante aux dés : après 20 lancers, une pièce défectueuse devient « un 5 ou un 6 ». Combien de lancers avant l'alarme ?\end{tcolorbox}
\vfill{\footnotesize\color{uimm}\textbf{Ressources de la base :} Tirer un lot · Jeu de secours.}
\end{document}